\par\Needspace{7\baselineskip}\bigskip
\begin{center}
\pdfbookmark[2]{Gauss an Olbers, Göttingen, 23. April 1814}{pallas-brief-18140423}
\textsc{Gauss} an \textsc{Olbers}, Göttingen, 23. April 1814.
\end{center}
\noindent\hspace*{\parindent}. . . . .\quad Wie genau die 9 Pallas-Oppositionen mit meinen Störungsrechnungen übereinstimmen, habe ich Ihnen bereits geschrieben. Indess sind mir doch die Differenzen noch zu gross. Die Störungen durch Saturn und Mars können sie vielleicht noch etwas herunterbringen, aber nach einem gemachten Überschlage doch nicht sehr viel. Es scheint freilich etwas bedenklich, auf Resultate Gewicht zu legen, welche zu finden für jede Opposition gegen 1000 Gleichungen numerisch berechnet werden mussten. Da indessen diese 9000 Rechnungen alle doppelt gerechnet sind, und ich in die Berechnung der Gleichungen selbst die grösste Sorgfalt gelegt habe, so halte ich es für kaum zweifelhaft, dass meine Jupitersmasse noch eine Verbesserung nöthig hat. Den wahren Werth dieser Correction zu finden, müssen nothwendig erst die Störungen durch Saturn und Mars abgestreift werden; jene hat \textsc{Nicolai} beinahe fertig, diese habe ich angefangen. Indess habe ich doch der Ungeduld nicht widerstehen können, zu sehen, wie viel sich (noch ohne Rücksicht der {\AstroSymbols ♄}- \SourceLeaf{429}und {\AstroSymbols ♂}-Störungen) jene Differenzen herunterbringen lassen, wenn die Jupitersmasse selbst als unbekannte Grösse behandelt wird. Mit grosser Überraschung fand ich, dass die Übereinstimmung dadurch \textls[180]{ganz ausserordentlich} gebessert wird. Sehen Sie hier die Längendifferenzen
\begin{center}
\begingroup\small\setlength{\tabcolsep}{3pt}
\begin{tabular}{@{}rrrrrrrrc@{}}
\multicolumn{9}{l}{mit der alten {\AstroSymbols ♃}-Masse}\\[2pt]
\(-0''{,}5\)&\(-21''{,}0\)&\(-13''{,}2\)&\(+19''{,}6\)&\(-25''{,}2\)&\(+16''{,}6\)&\(+4''{,}5\)&\(-27''{,}7\)&—\\[4pt]
\multicolumn{9}{l}{mit der verbesserten}\\[2pt]
\(+3{,}5\)&\(-3{,}1\)&\(+2{,}3\)&\(-0{,}5\)&\(+1{,}3\)&\(-5{,}9\)&\(+1{,}9\)&\(-7{,}1\)&—
\end{tabular}\endgroup
\end{center}
\begin{center}
\begingroup\small\setlength{\tabcolsep}{2pt}
\begin{tabular}{@{}cccccccr@{}}
mittlerer&Fehler&ohne&Correction&der&Masse&. . . . . .&\(17''{,}3\)\\
»&»&mit&»&»&»&. . . . . .&\(4{,}3\).
\end{tabular}\endgroup
\end{center}
\par Ich nenne mittlern Fehler die Quadratwurzel aus dem mittlern Werth des Quadrats des Fehlers, d.\,i.\ aus der Summe der Quadrate mit der Anzahl dividirt. Die Summe der Quadrate ist
\begin{center}
\begin{tabular}{@{}lcr@{}}
ohne Correction&. . . . . .&2684\\
mit Correction&. . . . . .&167.
\end{tabular}
\end{center}
\par Auch die Breiten, die übrigens schon ohne die Correction recht gut harmonirten, stimmen mit derselben noch etwas besser unter einander:
\begin{center}
\begin{tabular}{@{}lcccr@{}}
[ohne Correction]&mittlerer&Fehler&. . . . . .&\(5''{,}5\)\\
{}[mit Correction]&»&»&. . . . . .&\(4{,}0\).
\end{tabular}
\end{center}
\par Die Verbesserung der Masse ist \(\frac{1}{43}\) der alten und zwar Vergrösserung. Freilich ist dies Resultat noch precär, aber ich vermuthe, dass sie doch, nachdem {\AstroSymbols ♄} und {\AstroSymbols ♂} zugezogen sein werden, nicht kleiner als \(\frac{1}{60}\) sein wird. Wie ich höre, hat \textsc{Laplace} selbst in der IV. Ausgabe der Expos. (welche ich noch nicht gesehen habe) die alte Masse nach den Saturnbeobachtungen etwas vergrössert, obwohl nur sehr wenig, ich glaube die Pallas (zumal wenn erst noch mehrere Oppositionen benutzt werden können) muss hier zuverlässigere Resultate geben. Ich wünsche übrigens, dass obige Mittheilung noch unter uns bleibe. Die alte Bestimmung gründete sich bloss auf den 4. Jupiterstrabanten, dessen Abstand \textsc{Laplace} \(8'16''\) setzte; \textsc{Triesnecker} fand \(1\frac{1}{2}''\) mehr, und man müsste noch \(1\frac{1}{2}''\) mehr nehmen, um meine Masse zu finden.
\par
